
\documentclass[11pt]{article}
\usepackage[margin=1in]{geometry}

% Core packages
\usepackage{amsmath,amssymb}
\usepackage{tikz-cd}
\usepackage{multicol}
\usepackage{xurl}
\usepackage[hidelinks]{hyperref}
\usepackage{booktabs}
\usepackage{tabularx}
\usepackage{graphicx}
\usepackage{caption}
\captionsetup{font=footnotesize,labelfont=bf}

% Paragraphs
\setlength{\parindent}{0pt}
\setlength{\parskip}{1\baselineskip}
\setlength{\emergencystretch}{3em}
\hbadness=10000
\hfuzz=30pt
\vfuzz=20pt

\title{Benchmarking General and Medical Small Language Models for Spanish Obstetric Question Answering: The Impact of Fine-Tuning}
\author{Jhon Hander Mejia, Nicolas Hoyos Giraldo, Kevin D. Trujillo, Daniel Betancur Vasquez\\
Institucion Universitaria de Envigado\\
Facultad de Ingenieria}
\date{}

\begin{document}
\raggedbottom
\maketitle

\begin{abstract}
\textbf{Background:} Small Language Models (SLMs) have emerged as an efficient alternative to large language models for clinical question answering due to their lower computational requirements and ease of deployment. Recent studies have shown that domain adaptation through fine-tuning and retrieval-augmented generation (RAG) can substantially improve medical reasoning and factual accuracy. However, current evidence has primarily focused on general medical benchmarks, with limited evaluation in obstetric question answering, particularly in Spanish. Moreover, it remains unclear whether medically pretrained SLMs continue to benefit from additional domain-specific fine-tuning or whether repeated specialization may reduce their generalization capabilities.

\textbf{Objective:} This study benchmarks general-purpose and medically pretrained SLMs to evaluate the impact of domain-specific fine-tuning on Spanish obstetric question answering. The objective is to determine whether prior medical pretraining provides additional benefits after adaptation to a specialized obstetric corpus.

\textbf{Methods:} Two open-weight SLMs were evaluated: a general-purpose Gemma model and MedGemma, a medically pretrained counterpart. Both models were fine-tuned using the MaternaQA\-es dataset, a Spanish question-answering corpus focused on obstetric care. Performance was evaluated before and after fine-tuning using multiple automatic metrics assessing answer quality, semantic similarity, factual consistency, and clinical relevance. To further assess model robustness, retrieval-augmented generation (RAG) configurations were incorporated during inference as an auxiliary evaluation strategy rather than as the primary contribution of the study.
\end{abstract}

\newpage

\section{Introduction}
\begin{multicols}{2}
\raggedright

The rapid adoption of Large Language Models (LLMs) has transformed the development of clinical decision-support systems and medical question-answering (QA) applications. Nevertheless, their computational requirements limit deployment in many healthcare environments, motivating increasing interest in Small Language Models (SLMs), which offer lower inference costs, reduced hardware requirements, and easier local deployment while maintaining competitive performance in domain-specific tasks \cite{Pingua2025MedicalLLMs}. In medical applications, adapting SLMs to specialized knowledge has become a key research direction, primarily through supervised fine-tuning, allowing compact models to acquire domain-specific knowledge while preserving computational efficiency.

Recent studies have shown that domain-specific fine-tuning substantially improves biomedical question answering across a wide range of language models, leading to higher factual consistency and improved clinical relevance compared with general-purpose models \cite{Bora2025BiomedicalQA}. Nevertheless, most evaluations have been conducted using general biomedical datasets, English-language corpora, or heterogeneous clinical specialties, leaving limited evidence regarding highly specialized domains such as Spanish obstetric question answering \cite{Arzideh2026ClinicalEmbedding}. Reviews focusing on maternal healthcare further emphasize the lack of standardized benchmarks and clinically validated evaluation protocols for obstetric language models, highlighting the need for specialized resources that support reliable, evidence-based clinical question answering while reducing hallucinations \cite{BioMedInformatics2025MaternalRAG,Psilopatis2026ObstetricLLM,PachecoRomero2023IA}.

Another important but largely unexplored question concerns the value of medical pretraining after domain adaptation. Medical language models are expected to encode biomedical terminology and clinical reasoning acquired during pretraining. However, once these models undergo additional fine-tuning on a highly specialized downstream task, it remains unclear whether their previous medical knowledge continues providing measurable advantages over general-purpose SLMs adapted with the same dataset. Existing studies have rarely addressed this question under controlled experimental settings, particularly in maternal healthcare.

An additional unanswered question concerns the role of medical pretraining after domain adaptation. While medically pretrained SLMs are expected to possess stronger biomedical knowledge, it remains unclear whether additional fine-tuning on a highly specialized obstetric corpus consistently improves performance or whether repeated specialization introduces diminishing returns or reduced generalization. Despite its practical importance for developing lightweight clinical assistants, this question has received little empirical attention in the literature.

To investigate this question, we employ MaternaQA-es, this study presents a systematic benchmark of two open-weight SLMs representing different initialization strategies: a general-purpose model (Gemma 4) and a medically pretrained model (MedGemma). Both models are evaluated before and after fine-tuning using MaternaQA\-es, a publicly available Spanish obstetric question-answering dataset comprising 5,727 synthetic QA pairs generated from 57 clinically curated source documents through a reproducible pipeline including document curation, segmentation, quality control, and grounding validation. Fine-tuning is performed under an identical QLoRA configuration implemented in the MaternaCare-ES framework, enabling a controlled comparison between baseline and adapted models using standardized automatic evaluation metrics.

Specifically, this work makes three contributions. First, it provides the first benchmark comparing a general-purpose and a medically pretrained SLM for Spanish obstetric question answering. Second, it evaluates the impact of domain-specific fine-tuning under identical training conditions, enabling an unbiased comparison between both initialization strategies. Finally, the resulting fine-tuned models establish the foundation for a subsequent retrieval-augmented evaluation aimed at analyzing the interaction between specialized language models and external clinical knowledge.

\end{multicols}


\enlargethispage{2\baselineskip}

\begin{multicols}{2}
\raggedright
\section{Methodology}

\subsection{Study Design and Research Questions}

The study uses a controlled, offline factorial benchmark. The experimental factors are model family, parameter-adaptation state, and retrieval strategy. Model family has two levels: Google Gemma 4 E2B instruction-tuned (`google/gemma-4-E2B-it`) and Google MedGemma 1.5 4B instruction-tuned (`google/medgemma-1.5-4b-it`). Parameter adaptation has two levels: base weights without an adapter and grounded QLoRA. Retrieval strategy has three levels: the no-retrieval control, whose exact implementation label is `no\_rag`, canonical Hybrid, and canonical independent HyDE. Their Cartesian product yields $2 \times 2 \times 3 = 12$ conditions.

The benchmark workflow combines these factors with two evaluation scenarios and preserves the boundary between retained fine-tuning outputs and pending complete RAG results. The unit of analysis is an item-level paired `qa\_id`. Every condition in the complete matrix is intended to answer the same 328 held-out MaternaQA-es test items, allowing within-item contrasts rather than comparisons between unrelated samples. The primary estimand for QLoRA is the paired difference between QLoRA and base outputs within a fixed model family and retrieval strategy. Retrieval contrasts compare strategies within a fixed family and adaptation state. The interaction estimand is the change in the QLoRA effect across retrieval strategies. For example, with $H$ denoting Hybrid,
\[
  \Delta_{\mathrm{int}} =
  (Y_{\mathrm{QLoRA},H}-Y_{\mathrm{base},H})
  -(Y_{\mathrm{QLoRA},N}-Y_{\mathrm{base},N}).
\]
Analogous contrasts are defined for HyDE and for the two model families. Here, $N$ denotes the no-retrieval control (`no\_rag`).

\subsection{Dataset, Provenance, and Split Integrity}

MaternaQA-es is a Spanish obstetric question-answering resource containing 5,727 question-answer pairs derived from clinically curated source PDFs. Its QA publication splits explicitly contain 52 source PDFs in train, 2 in validation, and 3 in test (57 PDF representations total). The documented item split contains 5,093 training items, 306 validation items, and 328 held-out test items. Test source PDFs do not overlap the train or validation source PDFs. The test set is identified by `qa\_id`, which is retained as the pairing key for all conditions. Dataset construction includes document curation, PDF or text extraction, cleaning, segmentation, synthetic question-answer generation, grounding checks, and quality control \cite{MaternaQAes2025}.

The primary grounded adaptation uses the `sft\_grounded` representation: the source context associated with an item is concatenated with the question. `sft\_closed\_book` is a separate question-only representation and is not used as a primary ablation in the retained comparison. The distinction matters because the word \textit{base} describes whether an adapter is loaded, whereas \textit{closed-book} describes whether source context is supplied.

Split integrity is interpreted at two levels. The held-out 328-item test set is fixed and is reused across conditions. The retrieval corpus contains 2,268 document chunks, including the 108 source chunks referenced by the 328 test questions. This access is appropriate for evaluating evidence-grounded answering when the available obstetric source corpus can be searched. It is not a strict unseen-document generalization test. Retrieval access to those chunks must not be described as train/test leakage of model parameters: it concerns inference-time evidence availability, whereas QLoRA supervision concerns the examples and parameters used during adaptation.

\subsection{Protocol Evaluation Scenarios}

Two protocol scenarios define the evaluation inputs. The first is the \textit{clinician-authored obstetric challenge set}: a ten-item set prepared by obstetricians from Universidad Pontificia Bolivariana and represented by `datasets/sample10.jsonl` (the `DATA\_GT` source). Its records do not provide document-level reference chunk IDs. The second is the \textit{complete held-out MaternaQA-es test benchmark}: the canonical 328-item test split with QA metadata and 108 reference chunks associated with the test questions. The challenge set supports answer-quality evaluation but does not provide gold document-chunk identifiers for document-level retrieval recall; the complete test benchmark supports both response metrics and reference-chunk availability checks. Both are protocol scenarios, not claims that all 12 factorial cells have been executed; the retained numerical evidence and the pending RAG matrix remain distinguished below.

\subsection{Model Families and Experimental Conditions}

The benchmark compares Google Gemma 4 E2B instruction-tuned (`google/gemma-4-E2B-it`) with Google MedGemma 1.5 4B instruction-tuned (`google/medgemma-1.5-4b-it`). For each family, the base condition loads the corresponding published model without an adapter. The grounded QLoRA conditions use the adapter IDs `iue-edu/MaternaCare-ES-gemma4-qlora` and `iue-edu/MaternaCare-ES-medgemma-qlora`, respectively. The answer model, decoding settings, prompt structure, and test items are shared across retrieval strategies within a family and adaptation state.

The four currently retained model-output summaries are a grounded subset of the factorial design: Gemma 4 base, Gemma 4 QLoRA, MedGemma 1.5 base, and MedGemma 1.5 QLoRA. Each retained evaluation contains 328 outputs and uses `sft\_grounded` source context plus question. These files are not complete RAG results and do not represent the full 12-cell matrix. The numerical RAG benchmark requires corresponding run artifacts for all relevant conditions before results are reported.

\subsection{Grounded QLoRA Fine-Tuning}

QLoRA is the parameter-adaptation method in the benchmark. It combines quantized base-model weights with trainable low-rank adapter matrices, reducing the memory required for supervised adaptation while preserving the base model as a separately identifiable initialization \cite{Dettmers2023QLoRA}. RAFT is a related approach to adapting language models for domain-specific RAG, but it is not an additional condition in this benchmark \cite{Zhang2024RAFT}. Both model families use the same MaternaQA-es training and validation partitions, the same grounded input convention, and the same declared training configuration. The comparison is therefore intended to vary model initialization and adapter state, not the data or optimization protocol.

`sft\_grounded` training examples use the dataset's conversational `Contexto fuente` representation, in which source context is presented with the question. This supervised signal differs from retrieval at inference: the context is pre-associated with each training or evaluation example, whereas RAG selects contexts from the 2,268-chunk corpus for each test question. The adapter is loaded only in QLoRA conditions. Base conditions use the same answer model without that adapter and are not interpreted as closed-book conditions when context is otherwise present.

\par\medskip
\noindent\begin{minipage}{\linewidth}
    \centering
    \includegraphics[width=\linewidth]{../../assets/img/pdf/qlora-fine-tuning.pdf}
    \captionof{figure}{Grounded QLoRA fine-tuning pipeline. Source context is presented with each question; quantized base-model weights remain frozen while trainable low-rank adapter parameters are optimized. The resulting adapter is loaded only for grounded QLoRA conditions.}
    \label{fig:qlora-pipeline}
\end{minipage}
\par\medskip
\subsection{Retrieval-Augmented Generation Protocol}

RAG is the second core method. The protocol is specified independently of the retained fine-tuning summaries so that retrieval strategy is a reproducible experimental factor rather than an informal robustness check. Figure~\ref{fig:retrieval-protocol} distinguishes the control and retrieval paths and clarifies that hypothetical text is not supplied as corpus evidence.

\par\medskip
\noindent\begin{minipage}{\linewidth}
    \centering
    \includegraphics[width=\linewidth]{../../assets/img/pdf/all-rag-architecture.pdf}
    \captionof{figure}{RAG architecture variants specified for the benchmark. The no-retrieval control, Hybrid retrieval, and independent HyDE differ in how evidence is prepared before the answer model receives the question and available context. Hypothetical text is query expansion only; retrieved corpus chunks form the evidence context.}
    \label{fig:retrieval-protocol}
\end{minipage}
\par\medskip
\subsubsection{Corpus provenance and chunk construction}

The retrieval corpus contains 2,268 chunks from the curated obstetric source documents. The source pipeline is: PDF extraction and cleaning, chunk construction, deduplication, and metadata assignment. Source pages are grouped by PDF and section, and paragraphs are accumulated within each group. The implementation estimates tokens as $\lceil 1.35 \times \text{word\_count}\rceil$ and targets 500--1200 estimated tokens. When a chunk exceeds the maximum and is flushed, the implementation computes the overlap with the model tokenizer through `tail\_tokens`, retaining the last `overlap\_tokens` real tokenizer tokens (default 100). The 500--1200 window and this overlap are operational design choices intended to balance answerable context, paragraph and section-boundary preservation, and prompt budget; they were not clinically optimized values. Forced flushes can produce final chunks shorter than the target minimum. Chunks are retained when they contain at least 180 estimated tokens and have `clinical\_score` $\geq 5$. Metadata include a stable chunk identifier and source-document provenance. The same corpus is used for dense, sparse, and fused retrieval.

The construction signal `clinical\_score` is defined from normalized text, configured clinical terms, configured action terms, word count, and reference-line density. Let $\operatorname{norm}(x)$ apply Unicode NFKD decomposition, remove combining marks, lowercase, replace non-ASCII letters, digits, and whitespace with spaces, and collapse whitespace; let $T$ be the clinical-term set, $A$ the action-term set, $W$ the Unicode word count, and $R$ the ratio of reference-matching nonempty lines to all nonempty lines. With $\mathbf{1}(\cdot)$ denoting an indicator function, the implementation is:
\[
\begin{aligned}
S(x)&=\max\{0,L(x)+A(x)+B(W)-P(R)\},\\
L(x)&=\sum_{t\in T:\,t\subseteq\operatorname{norm}(x)}
      [1+\mathbf{1}(t\text{ contains a space})],\\
A(x)&=\sum_{a\in A:\,a\subseteq\operatorname{norm}(x)}1,\\
B(W)&=2\mathbf{1}(W\geq120)+2\mathbf{1}(W\geq350),\\
P(R)&=3\mathbf{1}(R\geq0.35).
\end{aligned}
\]
Here $S(x)$ is the reported `clinical\_score`; each configured term contributes at most once when it occurs as a substring. This heuristic is a construction and filtering signal, not clinical validation, and runtime RAG does not recompute it. The corpus filters retain the threshold `clinical\_score` $\geq 5$ and the minimum of 180 estimated tokens.

\par\medskip
\noindent\begin{minipage}{\linewidth}
    \centering
    \includegraphics[width=\linewidth]{../../assets/img/pdf/corpus-retrieval.pdf}
    \captionof{figure}{MaternaQA-es corpus construction and retrieval resource. Curated source documents are extracted, segmented, filtered, assigned provenance metadata, and retained as the 2,268-chunk corpus used by dense, sparse, and fused retrieval.}
    \label{fig:corpus-pipeline}
\end{minipage}
\par\medskip
\subsubsection{Dense and sparse retrieval}

Dense retrieval uses `BAAI/bge-m3`. The SentenceTransformer implementation produces normalized dense vectors, explicitly runs on CPU, and uses `batch\_size=32`. The persistent local vector table is LanceDB. Because the benchmark creates no vector index, LanceDB search is exhaustive exact cosine search, not approximate nearest-neighbor (ANN) retrieval. ANN would require an approximate vector index; it is not used in this benchmark. The default dense retrieval depth is $k=5$.

Sparse retrieval is Okapi BM25 implemented with `rank\_bm25`/`BM25Okapi`. It is lexical retrieval over Unicode lowercase \texttt{\textbackslash w+} tokens. BM25 supplies sparse term-weight scores and ranked chunk identifiers; this description does not imply a sparse-vector index. Tokenization and preprocessing are held fixed across questions within a run. The retrieval outputs are recorded before answer generation so that retrieval quality and latency can be analyzed separately from response generation.

\subsubsection{Canonical Hybrid retrieval}

Hybrid retrieval fuses dense and BM25 rankings. For a requested final depth $k$, the candidate depth is
\[
  k_{\mathrm{candidate}}=\min\left(|C|,\max(4k,50)\right),
\]
where $C$ is the 2,268-chunk corpus. Each retriever contributes its ranked candidates, and the union is fused by deterministic reciprocal rank fusion (RRF) with `rank\_constant=60` \cite{Cormack2009RRF}. Ties are broken by stable chunk ID. The final top-$k$ chunks are formatted in descending fused rank and supplied to the answer model.

\subsubsection{Canonical independent HyDE retrieval}

HyDE is an independent retrieval strategy. A dedicated hypothetical-document generator produces a Spanish hypothetical clinical-document passage from the question; this generator is separate from the answer model and is not counted as the answer model's retrieved evidence. The canonical OpenAI provider is `gpt-5-mini-2025-08-07`. An optional Hugging Face provider is `Qwen/Qwen2.5-1.5B-Instruct`. The prompt requests a Spanish hypothetical clinical-document passage relevant to answering the obstetric question, rather than a final patient-facing answer.

The default HyDE generation limit is 256 new tokens. Sampling is disabled, so generation is deterministic under the declared provider and model settings. The cache path is separated by dataset mode, provider, and generator model. Row reuse additionally checks `qa\_id` and a hash of the question; generation, prompt, and provider settings are included in the cache identity. This behavior does not assert full dataset-file hashing for the cache. Dense retrieval then embeds the hypothetical text with the pinned BGE-M3 encoder and searches the corpus by exact cosine similarity. The resulting chunks, not the hypothetical text itself, are passed as evidence to the answer model.

\subsubsection{No-retrieval control}

The no-retrieval control is the condition whose exact implementation label is `no\_rag`. It returns no retrieved contexts and omits the retrieved-context block from the answer prompt, leaving only the answer instruction and question. Context-dependent retrieval metrics are therefore not defined for `no\_rag`; they are not assigned a value of zero.

\subsection{Inference and Prompt Controls}

All answer prompts are written in Spanish. The prompt contains an instruction to answer the obstetric question, a clearly delimited context block, and the question itself. For RAG, the prompt contains the retrieved texts with positional labels such as [1] and [2]. Provenance is retained in benchmark records and metadata, not asserted as content of the answer prompt. For the no-retrieval control (`no\_rag`), the context block is empty. The answer model receives the same question formatting, chat template, and generation settings across retrieval strategies. The separate prompt conventions make retrieved context distinguishable from the pre-associated context used by `sft\_grounded` fine-tuning. This is an intentional protocol distinction: `sft\_grounded` training examples use the dataset's conversational `Contexto fuente` representation, whereas RAG inference uses a separate single-user prompt with `Contexto recuperado` and positional chunk labels. The resulting prompt distribution shift is a potential QLoRA$\times$RAG interaction and transportability consideration, not an undocumented mismatch.

Inference uses the model's chat template and greedy decoding by default, with a maximum of 512 newly generated tokens. The configured end-of-sequence and padding identifiers are used, and only newly generated tokens are decoded and evaluated. Sampling is disabled unless a run configuration explicitly records a different setting. QLoRA conditions load the declared adapter on the corresponding base model; base conditions do not load an adapter. Retrieval generation and answer generation are separate stages, and HyDE's dedicated generator is not substituted for the answer model.

\subsection{Evaluation Framework and Metrics}

Two evaluator pathways are distinguished. The legacy retained evaluator summarizes the four grounded outputs that already exist. The current RAG evaluator is the pathway defined for the 12-condition matrix and includes retrieval-specific metrics and operational telemetry. These measurement domains are separated before paired contrasts and interaction analysis. In both pathways, RAGAS 0.4.3 is an external automatic evaluation layer: it is neither the retriever nor a clinical expert. Metrics and their application belong in Methods; numerical scores belong in Results.

\subsubsection{Retained fine-tuning-output evaluator}

The legacy retained evaluator uses RAGAS version 0.4.3 \cite{Es2023Ragas} with `gpt-5.4-mini` as the evaluator language model, `text-embedding-3-small` for embedding-based evaluation, and a maximum completion length of 2,048 tokens. Four metrics were retained: Faithfulness, Answer Relevancy, Answer Correctness, and Semantic Similarity. In the order Gemma 4 base, Gemma 4 QLoRA, MedGemma 1.5 base, and MedGemma 1.5 QLoRA, the retained summaries contain error counts of 1, 1, 5, and 6, respectively. These counts describe evaluator-output records in the retained summaries; they are not RAG benchmark scores and do not complete the factorial matrix.

\subsubsection{RAG benchmark evaluator}

The RAG benchmark evaluator uses RAGAS version 0.4.3 \cite{Es2023Ragas} with six metrics: Context Precision, Context Recall, Faithfulness, Answer Relevancy, Answer Correctness, and Semantic Similarity. The current command-line evaluator has `text-embedding-3-large` as its embedding default. This evaluator embedding is distinct from the BGE-M3 encoder used for retrieval and must not be conflated with the retrieval model.

Context Precision measures whether relevant retrieved chunks are ranked ahead of irrelevant chunks. Context Recall measures how much of the reference-supporting information is retrieved. These are retrieval-quality metrics: they assess evidence selection and ordering rather than the prose of the answer. Faithfulness measures whether claims in the answer are supported by the supplied context. Answer Relevancy measures whether the response addresses the question. Answer Correctness compares the response with the reference answer for factual and semantic agreement. Semantic Similarity measures semantic alignment between the generated and reference answers. The latter four primarily characterize response quality, although Faithfulness is specifically conditioned on retrieved evidence.

The evaluator also records operational analysis variables: input, output, and total token counts; retrieval, generation, and end-to-end latency; and output tokens per second. With $N_{\mathrm{out}}$ denoting generated output tokens and $T_{\mathrm{gen}}$ the measured generation-stage latency,
    \[
      \mathrm{TPS}_{\mathrm{out}}=\frac{N_{\mathrm{out}}}{T_{\mathrm{gen}}}.
    \]
    This is an observed wall-clock output rate, not hardware throughput. Local Hugging Face timing includes prompt formatting and tokenization, generation, CUDA synchronization, and decoding. OpenAI HyDE timing includes request/response wall time. No library-native tokens-per-second metric is claimed, and these rates are not treated as cross-provider comparable. These variables therefore support resource-use and responsiveness descriptions only; they do not measure medical correctness. For the no-retrieval control (`no\_rag`), Context Precision, Context Recall, and Faithfulness are not applicable because no retrieved contexts are supplied. They are recorded as missing by design, not as zero-valued observations. The other response-quality metrics remain conceptually distinct from retrieval-quality metrics.

\subsection{Analysis Plan}

The primary analysis is descriptive and item-level. For each condition, the analysis will report the number of valid outputs, metric distributions, paired differences by `qa\_id`, and operational summaries. Descriptive means and contrasts use finite, metric-valid values. Evaluator failures are reported and not imputed. A paired contrast for a metric requires valid values for both paired outputs for that metric. QLoRA will be compared with base within each retrieval strategy and model family. Retrieval strategies will be contrasted within each adaptation state and model family. Model-family contrasts will compare Gemma 4 with MedGemma 1.5 within matched adaptation and retrieval conditions.

The interaction analysis will estimate whether the paired QLoRA difference changes with retrieval strategy and whether that change differs by model family. All contrasts will preserve the same `qa\_id` items wherever both outputs are available. Because the complete matrix and corresponding artifacts are not yet retained, inferential tests are not claimed as completed. If a complete rerun is performed, any future inferential analysis must freeze the exact aggregation and multiplicity procedure before execution; the paired structure, missingness, and uncertainty intervals will then be documented before inferential results are interpreted.

\subsection{Reproducibility, Ethics, and Threats to Validity}

Reproducibility requires recording the MaternaQA-es source revision, retrieval-corpus hash, complete configuration, prompts, model identifiers and revisions, QLoRA adapter identifiers, provider and model settings for HyDE, cache identity, evaluator version, and decoding settings. The pinned BGE-M3 revision and deterministic RRF tie-breaking are part of this record. Reporting follows the transparency goals of TRIPOD-LLM and Data Statements for NLP \cite{Gallifant2025TRIPODLLM,Bender2018DataStatements}.

The benchmark uses synthetic question-answer pairs and curated source documents. Synthetic QA can omit clinically relevant ambiguity, contain generation artifacts, or inadequately represent real patient questions. Automatic metrics are proxies and are not substitutes for clinician or expert assessment. The study does not validate clinical safety, diagnostic accuracy, patient outcomes, or deployment performance. No patient intervention is performed, and the benchmark should not be used as a clinical decision tool.

The retrieval corpus contains 2,268 chunks, including the 108 chunks associated with the held-out test questions. Consequently, RAG evaluates answering when an available evidence corpus contains the source material; it does not evaluate generalization to unseen documents. This retrieval-access limitation is distinct from parameter leakage and is reported explicitly. A further threat is retrieval/evaluation coupling: the evaluator may use a different language model and embedding model than the retriever, but automated judgments can still reflect evaluator preferences. The retained primary comparison has no completed closed-book ablation, because `sft\_closed\_book` is a separate question-only variant not included in the primary matrix. Finally, absent or incomplete run artifacts can bias any comparison; for that reason, the four retained grounded summaries are labeled as a subset rather than presented as completed RAG evidence.
\end{multicols}

\begin{table}[htbp]
\centering
\scriptsize
\caption{Specified factorial experiment matrix.}
\label{tab:factorial-matrix}
\begin{tabularx}{\textwidth}{c l l l X}
\toprule
\textbf{ID} & \textbf{Model family} & \textbf{Adaptation} & \textbf{Retrieval} & \textbf{Condition definition} \\
\midrule
C01 & Gemma 4 & Base & no\_rag & Base weights; no retrieved context \\
C02 & Gemma 4 & Base & Hybrid & Base weights; canonical Hybrid retrieval \\
C03 & Gemma 4 & Base & HyDE & Base weights; canonical independent HyDE retrieval \\
C04 & Gemma 4 & QLoRA & no\_rag & Grounded QLoRA adapter; no retrieved context \\
C05 & Gemma 4 & QLoRA & Hybrid & Grounded QLoRA adapter; canonical Hybrid retrieval \\
C06 & Gemma 4 & QLoRA & HyDE & Grounded QLoRA adapter; canonical independent HyDE retrieval \\
C07 & MedGemma 1.5 & Base & no\_rag & Base weights; no retrieved context \\
C08 & MedGemma 1.5 & Base & Hybrid & Base weights; canonical Hybrid retrieval \\
C09 & MedGemma 1.5 & Base & HyDE & Base weights; canonical independent HyDE retrieval \\
C10 & MedGemma 1.5 & QLoRA & no\_rag & Grounded QLoRA adapter; no retrieved context \\
C11 & MedGemma 1.5 & QLoRA & Hybrid & Grounded QLoRA adapter; canonical Hybrid retrieval \\
C12 & MedGemma 1.5 & QLoRA & HyDE & Grounded QLoRA adapter; canonical independent HyDE retrieval \\
\bottomrule
\end{tabularx}
\par\smallskip
\parbox{\textwidth}{\footnotesize\textit{Note.} The complete design is 2 model families $\times$ 2 adaptation states $\times$ 3 retrieval strategies. The four currently retained fine-tuning summaries correspond to the grounded subset Gemma 4 base/QLoRA and MedGemma 1.5 base/QLoRA, each with 328 outputs and `sft\_grounded` context. They are not completed RAG results. Numerical claims for the full matrix require the corresponding run artifacts.}
\end{table}

\enlargethispage{2\baselineskip}

\newpage

\begin{multicols}{2}
\raggedright
\begin{thebibliography}{99}
    \bibitem{Pingua2025MedicalLLMs} Pingua, B., Sahoo, A., Kandpal, M., Murmu, D., Rautaray, J., Barik, R. K., \& Saikia, M. J. (2025). Medical LLMs: Fine-Tuning vs. Retrieval-Augmented Generation. \textit{Bioengineering, 12}(7), 687. \url{https://doi.org/10.3390/bioengineering12070687}

    \bibitem{Bora2025BiomedicalQA} Bora, A., \& Cuay\'{a}huitl, H. (2025). Biomedical Literature Q\&A System Using Retrieval-Augmented Generation (RAG). \textit{arXiv}. \url{https://arxiv.org/abs/2509.05505}

    \bibitem{Arzideh2026ClinicalEmbedding} Arzideh, F., et al. (2026). Improving Retrieval Augmented Generation for Health Care by Fine-Tuning Clinical Embedding Models: Development and Evaluation Study. \textit{Journal of Medical Internet Research, 28}, e82997. \url{https://doi.org/10.2196/82997}

    \bibitem{BioMedInformatics2025MaternalRAG} Applications and Challenges of Retrieval-Augmented Generation (RAG) in Maternal Health: A Multi-Axial Review of the State of the Art in Biomedical QA with LLMs. (2025). \textit{BioMedInformatics, 7}(4), 148. \url{https://doi.org/10.3390/biomedinformatics7040148}

    \bibitem{Psilopatis2026ObstetricLLM} Psilopatis, I., et al. (2026). The role of large language models in the diagnosis and management of obstetric patients: A pilot feasibility study using simulated obstetric cases. \textit{Archives of Gynecology and Obstetrics}. \url{https://doi.org/10.1007/s00404-026-08423-1}

    \bibitem{Arslan2025ObGynEmergencies} Arslan, C., et al. (2025). Evaluation of ChatGPT-5 responses in obstetric and gynecological emergencies: Concordance, readability, and clinical reliability. \textit{BMC Emergency Medicine, 25}. \url{https://doi.org/10.1186/s12873-025-01396-5}

    \bibitem{Xu2026ObstetLLM} Xu, X., et al. (2026). Obstet-LLM: Large Language Model for Early Prediction of SGA-LGA Newborns. \textit{Journal of Computer Science and Technology}. \url{https://doi.org/10.1007/s11390-025-5020-0}

    \bibitem{PachecoRomero2023IA} Pacheco-Romero, J. (2023). Inteligencia artificial en la pr\'{a}ctica de la ginecobstetricia, la investigaci\'{o}n y redacci\'{o}n cient\'{i}fica. \textit{Revista Peruana de Ginecolog\'{i}a y Obstetricia, 69}(4). \url{https://www.scielo.org.pe/scielo.php?script=sci_arttext&pid=S2304-51322023000400001}

    \bibitem{MaternaQAes2025} Minciencias Maternas. (2025). \textit{MaternaQA-es} (Version 1.0) [Computer software]. GitHub. \url{https://github.com/minciencias-maternas/MaternaQA-es}

    \bibitem{Dettmers2023QLoRA} Dettmers, T., Pagnoni, A., Holtzman, A., \& Zettlemoyer, L. (2023). QLoRA: Efficient Finetuning of Quantized LLMs. \textit{Advances in Neural Information Processing Systems, 36}, 10088--10115. \url{https://doi.org/10.52202/075280-0441}

    \bibitem{Zhang2024RAFT} Zhang, T., Patil, S. G., Jain, N., Shen, S., Zaharia, M., Stoica, I., \& Gonzalez, J. E. (2024). RAFT: Adapting Language Model to Domain Specific RAG. \textit{arXiv preprint arXiv:2403.10131}. \url{https://doi.org/10.48550/arXiv.2403.10131}

    \bibitem{Lewis2020RAG} Lewis, P., et al. (2020). Retrieval-Augmented Generation for Knowledge-Intensive NLP Tasks. \textit{arXiv preprint arXiv:2005.11401}. \url{https://doi.org/10.48550/arXiv.2005.11401}

    \bibitem{Es2023Ragas} Es, S., James, J., Espinosa-Anke, L., \& Schockaert, S. (2023). Ragas: Automated Evaluation of Retrieval Augmented Generation. \textit{arXiv preprint arXiv:2309.15217}. \url{https://doi.org/10.48550/arXiv.2309.15217}

    \bibitem{Gao2023HyDE} Gao, L., Ma, X., Lin, J., \& Callan, J. (2023). Precise Zero-Shot Dense Retrieval without Relevance Labels. \textit{ACL 2023}, 1762--1777. \url{https://doi.org/10.18653/v1/2023.acl-long.99}

    \bibitem{Cormack2009RRF} Cormack, G. V., Clarke, C. L. A., \& B{\"u}ttcher, S. (2009). Reciprocal Rank Fusion Outperforms Condorcet and Individual Rank Learning Methods. \textit{SIGIR 2009}, 758--759. \url{https://doi.org/10.1145/1571941.1572114}

    \bibitem{Gallifant2025TRIPODLLM} Gallifant, J., et al. (2025). The TRIPOD-LLM Reporting Guideline for Studies Using Large Language Models. \textit{Nature Medicine}. \url{https://doi.org/10.1038/s41591-024-03425-5}

    \bibitem{Bender2018DataStatements} Bender, E. M., \& Friedman, B. (2018). Data Statements for Natural Language Processing: Toward Mitigating System Bias and Enabling Better Science. \textit{Transactions of the Association for Computational Linguistics, 6}, 587--604. \url{https://doi.org/10.1162/tacl_a_00041}
\end{thebibliography}
\end{multicols}

\end{document}
